%%%%%%%%%%%%%%%%%%%%%%%%%%%%%%%%%%%%%%%%%%%%%%%%%%%%%%%%%%%%%%%%%%%%%%%%%%%%%%%%
%2345678901234567890123456789012345678901234567890123456789012345678901234567890
%        1         2         3         4         5         6         7         8

\documentclass[letterpaper, 10 pt, conference]{IEEEtran}  % Comment this line out if you need a4paper

%\documentclass[a4paper, 10pt, conference]{ieeeconf}      % Use this line for a4 paper

\IEEEoverridecommandlockouts                              % This command is only needed if 
                                                          % you want to use the \thanks command

%\overrideIEEEmargins                                      % Needed to meet printer requirements.

%In case you encounter the following error:
%Error 1010 The PDF file may be corrupt (unable to open PDF file) OR
%Error 1000 An error occurred while parsing a contents stream. Unable to analyze the PDF file.
%This is a known problem with pdfLaTeX conversion filter. The file cannot be opened with acrobat reader
%Please use one of the alternatives below to circumvent this error by uncommenting one or the other
%\pdfobjcompresslevel=0
%\pdfminorversion=4


%The following packages can be found on http:\\www.ctan.org
\usepackage{graphics} % for pdf, bitmapped graphics files
\usepackage{epsfig} % for postscript graphics files
\usepackage{mathptmx} % assumes new font selection scheme installed
\usepackage{times} % assumes new font selection scheme installed
\usepackage{cite}
\usepackage{amsmath,amssymb,amsfonts}
\usepackage{algorithmic}
\usepackage{graphicx}
\usepackage{booktabs}
\usepackage{textcomp}
\usepackage{xcolor}

\title{\LARGE \bf
VoDOCK: Generalizable Last-Centimeter Robot Docking via Visual Goal Conditioning
}


\author{Albert Author$^{1}$ and Bernard D. Researcher$^{2}$% <-this % stops a space
\thanks{*This work was not supported by any organization}% <-this % stops a space
\thanks{$^{1}$Albert Author is with Faculty of Electrical Engineering, Mathematics and Computer Science,
        University of Twente, 7500 AE Enschede, The Netherlands
        {\tt\small albert.author@papercept.net}}%
\thanks{$^{2}$Bernard D. Researcher is with the Department of Electrical Engineering, Wright State University,
        Dayton, OH 45435, USA
        {\tt\small b.d.researcher@ieee.org}}%
}


\begin{document}



\maketitle
\thispagestyle{empty}
\pagestyle{empty}


%%%%%%%%%%%%%%%%%%%%%%%%%%%%%%%%%%%%%%%%%%%%%%%%%%%%%%%%%%%%%%%%%%%%%%%%%%%%%%%%
\begin{abstract}

Autonomous mobile robot (AMR) docking requires precise terminal alignment, yet conventional systems often rely on explicit localization, range sensing, maps, or dock-specific controllers. We propose the Vision-only DOCKing Policy (VoDOCK), a goal-conditioned policy that maps the visual relation between a current observation and a goal image directly to closed-loop base-velocity trajectories. VoDOCK uses vision as its only exteroceptive modality: it retains bidirectional relational tokens from a frozen pretrained model before its pose head, resamples them, and fuses them with proprioceptive wheel-velocity history to condition a diffusion policy without explicit relative-pose estimation. The policy is trained from 145 demonstrations (2.09~h) collected with one dock in one environment and evaluated, without fine-tuning, under background, dock-identity, and target-category shifts. Across the three held-out settings, VoDOCK achieves a mean success rate of \textbf{XX\%}, improving over the strongest baseline by \textbf{XX} percentage points, with aggregate median terminal errors of \textbf{XX~cm} and \textbf{XX}$^\circ$, a successful-trial median time to success of \textbf{XX~s}, and an online planning throughput of \textbf{XX.X plans/s}. These results indicate that pre-head relational features can transfer single-dock demonstrations to retargetable precision control specified by a goal image.

\end{abstract}

\begin{IEEEkeywords}
Vision-Based Navigation, Robot Docking, Goal-Conditioned Imitation Learning, Diffusion Policies
\end{IEEEkeywords}
\begin{figure*}[t!]
\centering
\includegraphics[width=\textwidth]{figures/arch_draft.pdf}
\caption{Training and inference paths of VoDOCK. During training, frozen ReLoc3R relational features are precomputed against each episode's final-frame goal and read from an HDF5 cache; the trainable condition network and DiT learn by noise-prediction score matching. At deployment, the supplied goal is encoded once, the same bidirectional relational streams are computed online from each synchronized observation window, and the EMA DiT is integrated from Gaussian noise by DPM-Solver++(2M). The denormalized trajectory is continuity-smoothed and a timestamped prefix is published; subsequent observations close the replanning loop.}
\label{fig:vodock_architecture}
\end{figure*}



%%%%%%%%%%%%%%%%%%%%%%%%%%%%%%%%%%%%%%%%%%%%%%%%%%%%%%%%%%%%%%%%%%%%%%%%%%%%%%%%
\section{Introduction}
\IEEEPARstart{F}{or} an AMR to operate for extended periods without human intervention, it must not only navigate to a charging station but also autonomously complete the final alignment required for physical engagement with the charging terminal. This final stage differs fundamentally from conventional navigation in its precision requirements. Position errors of only a few centimeters or heading errors of a few degrees can prevent successful terminal engagement, and the robot must be able to recover from such errors through closed-loop feedback during the approach.

Conventional precision docking systems often combine LiDAR or depth sensing, SLAM, prior maps, fiducial markers, explicit relative-pose estimation, and dock-specific controllers. Such modular pipelines can provide high accuracy in structured environments, but deployment to new docks or environments may require additional sensor calibration, map updates, or redesign of dock-specific control rules. Learning-based policies offer an alternative by directly learning the mapping of observations to actions. However, when trained with limited data, they may exploit visual shortcuts by associating a specific dock appearance or background with the desired behavior, rather than learning a reusable relationship between the current observation and the docking goal.

This work deliberately considers a restrictive training regime: every task-specific demonstration was collected with the same charging station in the same environment. We ask whether a policy trained from only approximately two hours of such single-dock data can perform centimeter-level docking and then transfer, without fine-tuning, across progressively larger distribution shifts: a new background, a new charging station, and visual targets outside the charging-station category. Addressing this question requires choosing an appropriate representation for docking control rather than simply selecting a stronger visual backbone. A single-image appearance embedding can effectively encode what is visible in a scene, whereas docking fundamentally requires information about how the current observation differs from the desired goal observation in terms of relative translation and orientation. We therefore formulate docking around a pairwise relational representation in which the current and goal images are jointly processed, rather than independently encoded and subsequently compared. This goal-image interface offers a practical path to scaling precise local approach and alignment: a new target can be specified by a goal image captured at its desired relative configuration instead of collecting a new demonstration set or designing another target-specific controller.

In our implementation, we use ReLoc3R~\cite{c4} as a frozen visual backbone because it satisfies this representational requirement. In particular, it provides spatially resolved bidirectional relational tokens in which the current observation attends to the goal, and the goal attends to the current observation. ReLoc3R is used only as a pretrained feature extractor within VoDOCK; we neither modify ReLoc3R nor claim improvements in its localization performance.

A central design choice is where to interface with the pretrained backbone. Rather than using its final rotation and translation estimates, VoDOCK truncates the network immediately before the pose head and retains the bidirectional dense relational tokens. Although collapsing the decoder features into a global pose estimate is appropriate for localization, this representation may be restrictive for precision control. Global aggregation can discard local geometric relationships around contact regions, edges, and partially occluded structures; a single pose estimate suppresses correspondence ambiguity and confidence information; and camera-pose regression and base-velocity trajectory generation optimize fundamentally different objectives. Moreover, errors in coordinate conventions, metric scale, or camera-to-body calibration can propagate directly into a controller when an explicit pose is used as input. Instead, VoDOCK allows a task-specific control network to directly select and transform relational information that is useful for docking.

The retained bidirectional dense tokens are compressed into a fixed-length latent representation using a learnable resampling module and fused at the token level with wheel-velocity history. A Diffusion Transformer then predicts a 2-s future base-velocity trajectory. During deployment, the policy repeatedly replans from the latest camera observation, yielding closed-loop control throughout the final approach. The system uses vision as its only exteroceptive modality and requires only a monocular RGB camera, proprioceptive wheel-velocity history, and a supplied goal image. It does not require LiDAR, depth sensing, SLAM, prior maps, fiducial markers, an explicit final-pose vector, or a separate appearance branch.

VoDOCK is trained from 145 successful demonstrations (2.09~h) collected with one dock in one environment. Without fine-tuning, the same checkpoint achieves a mean held-out success rate of \textbf{XX\%} under background, dock-identity, and target-category shifts, improving over the strongest baseline by \textbf{XX} percentage points.

Our contributions are twofold. First, we introduce a pre-head relational control interface that preserves bidirectional, spatially dense current--goal information for last-centimeter docking without explicit pose estimation. Second, we demonstrate real-robot transfer from single-dock demonstrations across background, dock-identity, and target-category shifts, supported by controlled appearance, final-pose, directional-stream, and token-capacity ablations.


\section{Related Work}

\subsection{Goal-Conditioned Visual Navigation}
Image-goal navigation extends a fixed policy $\pi(A\mid O)$ to a goal-conditioned policy $\pi(A\mid O,G)$ by providing an image captured at the target location as the goal observation. NoMaD~\cite{c1} unifies goal-conditioned navigation and goal-agnostic exploration within a single diffusion policy and demonstrates long-horizon navigation in previously unseen environments. However, conventional ImageNav benchmarks typically define success using meter-scale goal radii and do not directly evaluate the last-centimeter alignment required for physical engagement with a charging terminal.

AnyImageNav~\cite{c2} treats the goal image as a geometric query and recovers a 6-DoF goal pose through dense visual correspondences. It reports pose-recovery errors of 0.27~m/3.41$^\circ$ on Gibson and 0.21~m/1.23$^\circ$ on HM3D, substantially improving the precision of last-meter ImageNav. In contrast, our work does not address global exploration or explicit 6-DoF goal-pose recovery. Instead, we focus on the final approach regime, where the robot directly converts goal-relative visual relationships into velocity sequences for precise local alignment. Physical docking is the sole task used for policy training; alignment to held-out non-dock objects is evaluated only as a cross-category stress test of the learned goal-relative control behavior.

Aim My Robot~\cite{meng2025aim} addresses precision local navigation to arbitrary nearby objects and combines multimodal perception with precise action prediction, trained primarily on large-scale photorealistic simulation data for sim-to-real transfer. Like VoDOCK, it targets centimeter-level positioning and evaluates generalization to unseen objects. However, Aim My Robot predicts object-relative navigation actions for general downstream tasks, whereas VoDOCK studies physical last-centimeter docking from a goal image and learns direct closed-loop base-velocity trajectories from a small set of real-world demonstrations without explicit goal-pose supervision.

\subsection{Diffusion Policies for Closed-Loop Control}

Diffusion Policy~\cite{c3} formulates action generation as a conditional denoising process, enabling the modeling of multimodal action distributions and temporally coherent action sequences. When combined with receding-horizon execution, the policy can repeatedly replan from the latest observation and thereby operate in closed loop.

In VoDOCK, diffusion serves as an action generator that prevents multiple plausible approach and recovery trajectories from collapsing into a single averaged behavior. We do not regard diffusion itself as the primary source of generalization. Rather, transfer to a previously unseen dock depends critically on the goal-relative representation used to condition the diffusion model. Our focus is therefore on how relational visual information between the current and goal observations is exposed to the control policy.

\subsection{Pairwise Visual Representation as a Policy Component}

ReLoc3R~\cite{c4} is a relative camera-pose regression model pretrained on approximately eight million posed image pairs. We do not propose a new variant of ReLoc3R, nor do we claim its final pose-estimation performance as a contribution of this work. Within VoDOCK, ReLoc3R is used as a frozen backbone that jointly processes the current and goal images and provides spatially dense bidirectional relational tokens. The proposed policy is therefore conceptually separated from the particular backbone implementation and can, in principle, be paired with other pretrained models that expose an equivalent pairwise relational interface.

Conventional localization pipelines pass decoder features through a pose head and collapse them into a single rotation and translation estimate. VoDOCK instead operates before this final reduction. Rather than consuming the output of a localization module, the policy directly reads and adapts the pre-head representation for task-specific control. Learnable Perceiver queries~\cite{jaegle2021perceiver} select and compress local relational information, a fusion Transformer models interactions among temporal context, wheel-velocity history, and bidirectional visual tokens, and a Diffusion Transformer~\cite{peebles2023dit} converts the resulting representation into a closed-loop base-velocity trajectory. The key distinction is therefore not the availability of pretrained relational tokens themselves, but the use of the representation before it is compressed into an explicit pose as an information interface for precision control.

Accordingly, the approximately two hours of data used in this study should not be interpreted as the total data required to train the visual backbone from scratch. Rather, it corresponds to the amount of task-specific real-robot demonstration data used to learn the docking behavior of VoDOCK. By keeping the pretrained relational representation fixed and learning only the task-specific compression, multimodal fusion, and action-generation modules from a limited number of real-robot demonstrations, VoDOCK leverages large-scale visual pretraining while maintaining a low task-specific data requirement.


\section{Proposed Method}
\label{sec:method}

The proposed Vision-only DOCKing Policy (VoDOCK) is a closed-loop base-velocity policy that aligns the current observation with a desired target-relative observation specified by a goal image. Vision is the only exteroceptive modality; wheel-derived velocity history provides proprioceptive motion context. As illustrated in Fig.~\ref{fig:vodock_architecture}, VoDOCK consists of four stages: (i) current--goal relational feature extraction, (ii) task-adaptive token resampling, (iii) fusion of visual tokens and velocity history, and (iv) diffusion-based base-velocity trajectory generation. The trainable condition and action networks are shared exactly between training and deployment, whereas the surrounding execution paths differ: training reads precomputed relational features and performs score matching, while deployment computes those features online and performs reverse diffusion followed by robot-command postprocessing. The pretrained visual backbone remains frozen in both paths.

\subsection{Problem Formulation}
\label{subsec:problem_formulation}

At time $t$, the policy observation $O_t$ comprises an RGB observation history $\mathcal{I}_t$ containing five images and a 60-step wheel-derived velocity history $W_t$:

\begin{equation}
\begin{aligned}
    O_t=(\mathcal{I}_t,W_t),\\
    \mathcal{I}_t=\{I_{t-59},I_{t-47},I_{t-35},I_{t-23},I_{t-11}\},\\
    W_t=\{w_{t-j}\}_{j=0}^{59}.
\end{aligned}
\end{equation}

Camera images and wheel-derived velocities are recorded at 30~Hz. The velocity history contains 60 steps, corresponding to a 2.0-s observation buffer. The implementation selects visual indices $[0,12,24,36,48]$ from the oldest end of that buffer, giving the five observations shown above at 0.4-s intervals. Consequently, wheel history reaches time $t$, whereas the newest visual input is $I_{t-11}$; training and deployment use the same convention. A single goal image $G$, captured at the desired relative target configuration, specifies the control target. For a dock, this image depicts the desired terminal docking configuration; for a non-dock object, it depicts the desired approach and alignment configuration. Conditioned on $O_t$ and $G$, VoDOCK predicts $H_a=60$ future base-velocity commands, corresponding to a 2.0-s action horizon:

\begin{equation}
\begin{aligned}
A_t&=[a_t,\ldots,a_{t+H_a-1}]\in\mathbb{R}^{60\times2},\\
a_{t+h}&=(v_{t+h},\omega_{t+h}),\quad h=0,\ldots,H_a-1,\\
A_t&\sim p_\theta(A_t\mid O_t,G).
\end{aligned}
\end{equation}

During training, the final docked frame of each successful episode is used as $G$. The policy therefore learns to align the current observation with a goal observation rather than regress a fixed dock pose. At deployment, the policy can be retargeted by replacing only the goal image, which need not depict a charging station.

\subsection{Pre-Head Relational Visual Representation}
\label{subsec:relational_representation}

Each history image $I_t^{(i)}$ and the goal image $G$ are jointly processed as an image pair by a frozen ReLoc3R backbone~\cite{c4}. Its bidirectional cross-attention decoder produces two relational token streams associated with the two relative-transform directions: $Z_i^{c\rightarrow g}$ for current-to-goal and $Z_i^{g\rightarrow c}$ for goal-to-current:

\begin{equation}
\begin{aligned}
\left(Z_i^{c\rightarrow g},Z_i^{g\rightarrow c}\right)
    &=F_\phi\!\left(I_t^{(i)},G\right),\\
Z_i^{c\rightarrow g},Z_i^{g\rightarrow c}
    &\in\mathbb{R}^{196\times768}.
\end{aligned}
\end{equation}

where the backbone parameters $\phi$ remain fixed throughout policy training. This transform-based notation is used throughout the paper. The two streams encode the same image pair in opposite directions, retaining complementary correction and reference cues.

For policy training, encoder features are cached for every dataset frame and the final normalized decoder streams are precomputed by pairing each frame with its episode goal. The policy therefore reads $Z_i^{c\rightarrow g}$ and $Z_i^{g\rightarrow c}$ from an HDF5 sidecar without differentiating through $F_\phi$. At deployment, the goal image is encoded once during initialization; for every subsequent inference request, the five current images are batch-encoded and passed with the cached goal encoding through the same frozen bidirectional decoder. The live decoder outputs are quantized through the same FP16 tensor interface used by the training cache before entering the policy. Thus the semantic representation, shape, and downstream network are shared, while relational extraction is offline during training and online during inference.

VoDOCK discards the localization-specific pose head and does not use its final rotation or translation estimates. Although a global pose is appropriate for localization, the pose head compresses dense correspondences, local structures, partial occlusions, and correspondence ambiguity into a single point estimate. It can also propagate coordinate-frame, metric-scale, and camera-to-body calibration errors directly to a downstream controller. VoDOCK instead exposes the dense pre-head tokens as the policy interface and lets the docking action objective select the relational cues that are useful for control. This design does not modify ReLoc3R itself; it redesigns the interface between a pretrained relational representation and a task-specific control policy.

\subsection{Task-Adaptive Token Resampling and Fusion}
\label{subsec:token_fusion}

A linear projection first maps each directional token stream from 768 to $d=384$ dimensions. A Perceiver resampler~\cite{jaegle2021perceiver} with 16 learnable latent queries then compresses the 196 spatial tokens of each frame into a fixed-length representation:


\begin{equation}
\begin{aligned}
\widetilde{Z}_i^{c\rightarrow g}
    &=\mathcal{R}_{Q_c}\!\left(P_c Z_i^{c\rightarrow g}\right),\\
\widetilde{Z}_i^{g\rightarrow c}
    &=\mathcal{R}_{Q_g}\!\left(P_g Z_i^{g\rightarrow c}\right),\\
\widetilde{Z}_i^{(\cdot)}
    &\in\mathbb{R}^{16\times384}.
\end{aligned}
\end{equation}

Unlike spatial average pooling, learned resampling allows different latent queries to select distinct local and relational cues according to the docking objective. The five observation times and two relation directions produce 160 visual tokens. These tokens are concatenated with 60 wheel-velocity tokens and jointly updated by a four-layer, six-head self-attention Transformer:

\begin{equation}
    C_t=\mathcal{F}_\psi
    \left(
    \left[
    \{\widetilde{Z}_i^{c\rightarrow g}\}_{i=1}^{5};
    \{\widetilde{Z}_i^{g\rightarrow c}\}_{i=1}^{5};
    E_w(W_t)
    \right]
    \right)
    \in\mathbb{R}^{220\times384}.
\end{equation}


where $E_w$ embeds the wheel-velocity history into the token space and $\mathcal{F}_\psi$ denotes the multimodal fusion network. The resulting condition $C_t$ is passed to the action denoiser as a token sequence without global pooling.

\subsection{Goal-Conditioned Diffusion Action Policy}
\label{subsec:diffusion_policy}

VoDOCK models the conditional action distribution $p_\theta(A_t\mid O_t,G)$ using a diffusion model. During training, Gaussian noise is added to a clean action trajectory $A_t^0$ at diffusion step $k$:

\begin{equation}
A_t^k=\alpha_k A_t^0+\sigma_k\epsilon,
\qquad
\epsilon\sim\mathcal{N}(0,I).
\end{equation}

A 12-layer cross-attention Diffusion Transformer~\cite{peebles2023dit} allows the noisy action tokens to attend to the condition sequence $C_t$. The network minimizes the noise-prediction objective

\begin{equation}
    \mathcal{L}_{\mathrm{diff}}
=
\mathbb{E}_{A_t^0,k,\epsilon}
\left[
\left\|
\epsilon-\epsilon_\theta(A_t^k,k,C_t)
\right\|_2^2
\right].
\end{equation}

The diffusion formulation represents multiple feasible approach and recovery behaviors as temporally coherent base-velocity trajectories rather than collapsing them into a single averaged action. All modules downstream of the frozen relational backbone---including the projections, resamplers, wheel encoder, fusion Transformer, and action denoiser---are optimized end-to-end. This forward-noising objective and its gradient update are used only during training; deployment instead solves the learned reverse process without ground-truth future actions or gradient computation.

\subsection{Online Inference and Closed-Loop Execution}
\label{subsec:closed_loop}

At startup, the deployment process reconstructs the sensor encoders, fusion Transformer, and DiT from the configuration stored beside the checkpoint and loads both the online and exponential-moving-average (EMA) parameter sets. It also encodes the supplied goal image once with frozen ReLoc3R. Each inference request then reads a synchronized 60-frame RGB and wheel-velocity window, applies the same temporal indices and min--max statistics as training, computes the two live relational streams, and produces $C_t\in\mathbb{R}^{220\times384}$ with the trained condition network.

Sampling begins from a Gaussian action trajectory and uses the EMA model with DPM-Solver++(2M)~\cite{lu2022dpmsolverpp}:

\begin{equation}
    A_t^{K}\sim\mathcal{N}(0,I),\qquad
    \widehat{A}_t^{0}
    =\operatorname{DPM}_{\theta_{\mathrm{EMA}}}
    \!\left(A_t^{K},C_t;K=30\right).
\end{equation}

The 12-layer DiT is therefore evaluated repeatedly along 30 reverse-denoising steps, with all 60 action tokens jointly self-attending and cross-attending to $C_t$ at every layer. We use one trajectory sample per inference request; hence the configured medoid aggregation is the identity. The normalized result is mapped back to physical linear and angular velocity using the action statistics stored in the checkpoint. Model-weight EMA should not be confused with trajectory EMA: the deployed plugin executes the newly sampled plan rather than its trajectory-EMA output and subsequently applies a separate continuity filter. For action index $i$, this filter is

\begin{equation}
\begin{aligned}
\bar{a}_{t,i}
  &=\beta_i\widehat{a}_{t,i}^{0}
    +(1-\beta_i)\bar{a}_{t-1,i+1},\\
\beta_i
  &=\frac{1-\exp(-\gamma i)}
          {1-\exp[-\gamma(H_a-1)]},\qquad \gamma=0.2,
\end{aligned}
\end{equation}

where the final element of the shifted previous plan is repeated and the first plan is left unchanged. The model generates 60 commands, while the robot interface retains the first $K_e=32$ commands (1.07~s at 30~Hz). These velocities are integrated into cumulative forward-distance and heading increments and then packaged as a timestamped 32-step command prefix. With the evaluated runtime settings, this prefix is published every 0.1~s and replaced whenever a new inference result arrives. This fixed 10-Hz publication schedule is distinct from policy throughput: publication may reuse the latest plan, whereas throughput counts newly completed plans. A new synchronized observation window then updates the current--goal relation and closes the receding-horizon loop. If the required sensor window is incomplete or inference raises an exception, the deployment process emits a zero-velocity prefix. No LiDAR, depth, SLAM, prior map, explicit goal-pose estimate, or separate classical goal controller is used.


\section{Experimental Evaluation}
\label{sec:experiments}

We evaluate whether VoDOCK (i) learns repeatable docking from single-dock demonstrations, (ii) remains robust to background changes, (iii) transfers to a held-out charging station, and (iv) supports goal-conditioned alignment to objects outside the training category. All methods are evaluated from the same initial-condition distribution under identical success criteria.

\subsection{Setup and Evaluation Protocol}
\label{subsec:evaluation_protocol}

Tables~\ref{tab:training_configuration} and~\ref{tab:inference_configuration} summarize the task-specific training and deployment configurations. All 145 successful demonstrations were collected with one charging station in one environment. The dataset contains 225,465 frames sampled at 30~Hz, corresponding to 7,515.5~s (2.09~h) of real-world trajectories. An episode lasts 51.83~s on average and 48.50~s at the median. VoDOCK is trained for 20 epochs and 16,940 gradient updates without task-specific pose supervision. The held-out station, changed-background scenes, and non-dock targets contribute no task-specific training data. At test time, a single frozen checkpoint is retargeted only by changing the goal image. The deployment runtime can also load ablation checkpoints, but only the dense pre-head \texttt{r\_relfeat\_only} variant constitutes VoDOCK in this evaluation. Figure~\ref{fig:cross_dock_setup} shows the robot platform and the seen and held-out docking stations.

\begin{table}[!t]
\caption{Task-Specific Dataset and VoDOCK Training Configuration}
\label{tab:training_configuration}
\centering
\footnotesize
\renewcommand{\arraystretch}{1.05}
\begin{tabular*}{\columnwidth}{@{\extracolsep{\fill}}lr@{}}
\toprule
Item & Configuration \\
\midrule
Successful demonstrations & 145 \\
Training dock instances & 1 \\
Training environments & 1 \\
Total frames & 225,465 \\
Sampling frequency & 30~Hz \\
Total demonstration time & 7,515.5~s (2.09~h) \\
Episode duration & 51.83~s mean / 48.50~s median \\
Training epochs & 20 \\
Gradient updates & 16,940 \\
Batch size & 256 \\
Optimizer & AdamW \\
Learning rate / weight decay & $10^{-4}$ / $10^{-5}$ \\
EMA decay & 0.999 \\
Action normalization & Min--max \\
Task-specific pose supervision & None \\
\bottomrule
\end{tabular*}
\end{table}

\begin{table}[!t]
\caption{VoDOCK Real-Robot Inference Configuration}
\label{tab:inference_configuration}
\centering
\footnotesize
\renewcommand{\arraystretch}{1.05}
\begin{tabular*}{\columnwidth}{@{\extracolsep{\fill}}lr@{}}
\toprule
Item & Configuration \\
\midrule
Parameter set & Model-weight EMA \\
Diffusion solver & DPM-Solver++(2M) \\
Reverse steps / samples & 30 / 1 \\
Generated horizon & 60 steps (2.0~s) \\
Published prefix & 32 steps (1.07~s) \\
Trajectory EMA & Disabled (current plan) \\
Continuity coefficient $\gamma$ & 0.2 \\
Command publication interval & 0.1~s \\
Goal encoding & Once at initialization \\
Current--goal relations & Online at each inference \\
Buffer replacement & On new inference result \\
\bottomrule
\end{tabular*}
\end{table}

\begin{figure*}[t]
\centering
\begin{minipage}[b]{0.28\textwidth}
    \centering
    \includegraphics[width=\linewidth]{figures/amr1.png}
    \par\smallskip
    {\small\textbf{(a) AMR platform}}
\end{minipage}
\hfill
\begin{minipage}[b]{0.28\textwidth}
    \centering
    \includegraphics[width=\linewidth]{figures/seen_dock.png}
    \par\smallskip
    {\small\textbf{(b) Training: seen dock}}\\[-0.2em]
    {\footnotesize One environment}
\end{minipage}
\hfill
\begin{minipage}[b]{0.28\textwidth}
    \centering
    \includegraphics[width=\linewidth]{figures/unseen_dock.png}
    \par\smallskip
    {\small\textbf{(c) Test: held-out dock}}\\[-0.2em]
    {\footnotesize No fine-tuning}
\end{minipage}
\caption{Training hardware and evaluation settings. Task-specific demonstrations use only the seen dock in one environment. The same checkpoint is evaluated under background, dock-identity, and target-category shifts by changing only the goal image.}
\label{fig:cross_dock_setup}
\end{figure*}

\begin{figure*}[t]
\centering
\includegraphics[height=0.20\textheight,keepaspectratio]{figures/ex1_frontal_approach.png}
\hfill
\includegraphics[height=0.20\textheight,keepaspectratio]{figures/ex3_neardock_recv.png}
\hfill
\includegraphics[height=0.20\textheight,keepaspectratio]{figures/ex2_off_heading_test.png}
\caption{Docking-station evaluation configurations. (a) Frontal approaches at nominal 1~m (five configurations) and 2~m (seven configurations), with the angular locations shown. (b) Near-dock recovery from two lateral-offset poses, annotated as 10~cm forward and 20~cm to either side. (c) Two opposite initial headings at the same nominal 1-m centerline position. Triangles denote initial robot poses; the diagrams are schematic and not to scale.}
\label{fig:docking_scenarios}
\end{figure*}

Figure~\ref{fig:docking_scenarios} defines 16 docking-station initial configurations grouped into three complementary scenario families. The frontal-approach family varies both standoff distance and angular location while initializing the robot toward the dock. The near-dock family tests lateral correction close to the terminal configuration, and the off-heading family isolates heading recovery by keeping the 1-m centerline position fixed while reversing the initial viewing direction. These markers specify initial poses rather than executed trajectories. Every compared method uses the same configuration list and reset procedure within a given dock and environment condition.

We test the seen dock in the training background, the seen dock in a different background, a held-out dock, and held-out non-dock targets. For each docking-station condition, trials follow the predefined initial-pose cases in Fig.~\ref{fig:docking_scenarios}, with any repetitions balanced across methods. Dock success requires collision-free physical engagement within the time limit. For non-dock targets, success denotes terminal alignment within 3~cm and $5^\circ$ of the goal-image configuration. Terminal position and heading errors are measured using \textbf{XX measurement system or procedure}. Mean held-out success is the unweighted mean over the changed-background, held-out-dock, and held-out non-dock conditions.

In addition to success rate and terminal error, we report time to success (TTS) and online planning throughput. For a successful trial $j$, TTS is measured from the first nonzero robot command to physical engagement for a dock, or to the first satisfaction of the terminal tolerance for a non-dock goal. We report the median and interquartile range (IQR) over successful trials. Failed trials remain failures in the success-rate analysis and are treated as right-censored at the trial time limit rather than being assigned an artificial completion time.

Planning throughput counts complete predicted trajectories, not individual action samples or command publications. For $N$ sequential batch-one inference requests, let $t_i^{\mathrm{in}}$ denote the time immediately before the synchronized snapshot enters the policy and $t_i^{\mathrm{out}}$ the time at which its 32-step command prefix is ready. We compute

\begin{equation}
Q_{\mathrm{plan}}
=\frac{N}{\sum_{i=1}^{N}\left(t_i^{\mathrm{out}}-t_i^{\mathrm{in}}\right)}
\quad [\mathrm{plans\,s^{-1}}].
\end{equation}

The timed region includes online ReLoc3R relation extraction, condition encoding, 30-step reverse diffusion, denormalization, continuity filtering, and prefix construction. It excludes one-time model initialization and goal encoding, sensor waiting, interprocess-queue delay, and MQTT transport. All methods use one trajectory sample per request and identical solver settings unless stated otherwise. On the deployment \textbf{XX CPU and XX GPU}, we discard 10 warm-up requests and time 100 valid requests, synchronizing the GPU at both timing boundaries. We report $Q_{\mathrm{plan}}$ together with median and 95th-percentile end-to-end plan latency. Higher throughput and lower TTS are better.

\textbf{Baselines and controls:}
The \emph{DINO appearance-only} baseline independently encodes the current and goal images using frozen DINOv3 features~\cite{dino2025}, testing whether joint pairwise reasoning is necessary. Its five history encodings, one goal encoding, and velocity history produce 156 condition tokens. The \emph{final-pose interface} replaces the dense relational tokens with ReLoc3R's two 12-D rotation--translation outputs, testing whether explicit bidirectional relative poses are sufficient for direct control. Each pose stream is projected and resampled with 16 queries per frame, retaining VoDOCK's 220-token condition length. Two directional ablations isolate the relational streams. The \emph{current-to-goal} model retains only $Z^{c\rightarrow g}$, whereas the \emph{goal-to-current} model retains only $Z^{g\rightarrow c}$. Each directional model first uses the same 16 queries per frame as one VoDOCK stream, producing 140 condition tokens including velocity history. Its token-matched counterpart uses 32 queries per frame and restores the 220-token condition length of VoDOCK. Evaluating both configurations separates directional complementarity from reduced token capacity.

% TODO before reporting a controlled appearance comparison: retrain
% dino_goal_only with the same data split, batch size (256), optimizer,
% epochs, seed policy, action horizon, and diffusion stack as VoDOCK.

\subsection{Real-Robot Results}
\label{subsec:robot_results}

Table~\ref{tab:robot_results_en} reports docking or goal-alignment success, successful-trial TTS, and batch-one online planning throughput for every model, together with VoDOCK's terminal accuracy. Across the three held-out conditions, VoDOCK achieves a mean success rate of \textbf{XX\%}, improving over the strongest baseline by \textbf{XX} percentage points. Its aggregate median terminal position and heading errors are \textbf{XX~cm} and \textbf{XX}$^\circ$, respectively. VoDOCK reaches successful termination in a median \textbf{XX~s} [IQR: \textbf{XX--XX~s}] and produces \textbf{XX.X plans/s}, with median and 95th-percentile end-to-end planning latencies of \textbf{XX~ms} and \textbf{XX~ms}.

\begin{table*}[!t]
\caption{Real-robot docking and goal-alignment results. Panel A reports successful trials/total trials and batch-one online planning throughput; throughput counts complete 60-step plans per second and is independent of the 10-Hz command publication schedule. Panel B reports TTS in seconds as median [IQR] over successful trials; lower is better. Failed trials remain represented in Panel A and are right-censored at the time limit. Token-matched variants use 32 queries per frame to match VoDOCK's 220 condition tokens. The final rows of Panel A report VoDOCK's median terminal errors.}
\label{tab:robot_results_en}
\centering
\scriptsize
\textbf{Panel A: Success and online planning throughput}\\[0.25em]
\resizebox{\textwidth}{!}{%
\begin{tabular}{@{}lcccccc@{}}
\toprule
Model or metric & Tokens & Throughput (plans/s) $\uparrow$ & Seen dock, train bg. & Seen dock, new bg. & Held-out dock & Held-out non-dock goals \\
\midrule
DINO appearance only                         & 156 & XX.X & XX/XX & XX/XX & XX/XX & XX/XX \\
Final-pose interface                         & 220 & XX.X & XX/XX & XX/XX & XX/XX & XX/XX \\
Current-to-goal (16 queries)                 & 140 & XX.X & XX/XX & XX/XX & XX/XX & XX/XX \\
Current-to-goal (32 queries, token-matched)  & 220 & XX.X & XX/XX & XX/XX & XX/XX & XX/XX \\
Goal-to-current (16 queries)                 & 140 & XX.X & XX/XX & XX/XX & XX/XX & XX/XX \\
Goal-to-current (32 queries, token-matched)  & 220 & XX.X & XX/XX & XX/XX & XX/XX & XX/XX \\
\textbf{VoDOCK}                              & 220 & \textbf{XX.X} & \textbf{XX/XX} & \textbf{XX/XX} & \textbf{XX/XX} & \textbf{XX/XX} \\
\midrule
VoDOCK median position error (cm) $\downarrow$      & -- & -- & XX & XX & XX & XX \\
VoDOCK median heading error ($^\circ$) $\downarrow$ & -- & -- & XX & XX & XX & XX \\
\bottomrule
\end{tabular}%
}
\par\smallskip
\resizebox{\textwidth}{!}{%
\begin{tabular}{@{}lccccc@{}}
\toprule
\multicolumn{6}{c}{\textbf{Panel B: Time to success (s), median [IQR] over successful trials $\downarrow$}} \\
\midrule
Model & Tokens & Seen dock, train bg. & Seen dock, new bg. & Held-out dock & Held-out non-dock goals \\
\midrule
DINO appearance only                         & 156 & XX [XX--XX] & XX [XX--XX] & XX [XX--XX] & XX [XX--XX] \\
Final-pose interface                         & 220 & XX [XX--XX] & XX [XX--XX] & XX [XX--XX] & XX [XX--XX] \\
Current-to-goal (16 queries)                 & 140 & XX [XX--XX] & XX [XX--XX] & XX [XX--XX] & XX [XX--XX] \\
Current-to-goal (32 queries, token-matched)  & 220 & XX [XX--XX] & XX [XX--XX] & XX [XX--XX] & XX [XX--XX] \\
Goal-to-current (16 queries)                 & 140 & XX [XX--XX] & XX [XX--XX] & XX [XX--XX] & XX [XX--XX] \\
Goal-to-current (32 queries, token-matched)  & 220 & XX [XX--XX] & XX [XX--XX] & XX [XX--XX] & XX [XX--XX] \\
\textbf{VoDOCK}                              & 220 & \textbf{XX [XX--XX]} & \textbf{XX [XX--XX]} & \textbf{XX [XX--XX]} & \textbf{XX [XX--XX]} \\
\bottomrule
\end{tabular}%
}
\end{table*}

\textbf{Seen-dock performance:}
In the training background, VoDOCK succeeds in \textbf{XX/XX} trials, establishing that the proposed representation retains the precision required for physical engagement. When only the background changes, it succeeds in \textbf{XX/XX} trials, a change of \textbf{XX} percentage points, indicating that the policy does not rely solely on the training scene.

\textbf{Target transfer:}
Without fine-tuning, VoDOCK succeeds in \textbf{XX/XX} trials on the held-out dock and \textbf{XX/XX} trials on held-out non-dock goals. The latter condition is not interpreted as physical docking; it is a cross-category stress test of whether replacing the goal image can retarget the learned local alignment behavior.

\subsection{Ablation Study}
\label{subsec:ablation_analysis}

The ablations in Table~\ref{tab:robot_results_en} are organized around two questions. First, the DINO appearance-only and final-pose interfaces test whether independent appearance features or explicit bidirectional relative poses preserve enough information for docking control. Second, the current-to-goal and goal-to-current variants isolate the contribution of each relational direction. For each single-direction model, the 16-query setting retains the per-stream resampling capacity used by VoDOCK, whereas the 32-query token-matched setting restores the total condition length from 140 to 220 tokens. Comparing the two settings separates the effect of condition-token capacity from that of removing a relational direction. Panel B tests whether a representation changes the time required for successful closed-loop convergence, while the throughput column quantifies its online computational cost. Overall, VoDOCK improves mean held-out success by \textbf{XX} percentage points over the strongest trained baseline, with a successful-trial median TTS difference of \textbf{XX~s} and a throughput difference of \textbf{XX.X plans/s} relative to that baseline.

\textbf{Relational representation:}
VoDOCK improves over the DINO appearance-only and final-pose interfaces by \textbf{XX} and \textbf{XX} percentage points in mean held-out success, respectively. This result indicates that jointly processed, spatially dense pre-head features provide control-relevant information that is not preserved by independent appearance encoding or explicit pose compression.

\textbf{Directionality and token capacity:}
The 16-query current-to-goal and goal-to-current variants achieve \textbf{XX\%} and \textbf{XX\%} mean held-out success. Increasing the surviving stream to 32 queries changes these results to \textbf{XX\%} and \textbf{XX\%}. Comparing each pair isolates the effect of condition-token capacity, while the remaining gap to bidirectional VoDOCK measures the benefit of retaining both relational directions.

\section{Conclusion}
\label{sec:conclusion}

This work introduced VoDOCK, a goal-conditioned policy that preserves bidirectional relational tokens before the localization head and converts them, together with wheel-derived velocity history, into closed-loop base-velocity trajectories. Trained from 145 single-dock demonstrations totaling 2.09~h, VoDOCK achieves \textbf{XX\%} mean success across three held-out conditions, improves over the strongest baseline by \textbf{XX} percentage points, reaches success in a median \textbf{XX~s} [IQR: \textbf{XX--XX~s}], and produces \textbf{XX.X plans/s} on the deployment computer, without additional demonstrations, fine-tuning, explicit pose estimation, or range sensing.

The current study is limited to one robot and camera configuration, one training dock and environment, \textbf{XX} trials per condition, and a small set of held-out targets. Moreover, the 2.09~h budget excludes ReLoc3R pretraining and should be interpreted as task-specific transfer data rather than end-to-end visual training data. Broader hardware, lighting, connector, and object-category evaluations are required to characterize rare failures and the limits of goal-image retargeting.


\begin{thebibliography}{99}

\bibitem{c4} S. Dong, S. Wang, S. Liu, L. Cai, Q. Fan, J. Kannala, and Y. Yang, ``ReLoc3R: Large-Scale Training of Relative Camera Pose Regression for Generalizable, Fast, and Accurate Visual Localization,'' in \emph{Proc. IEEE/CVF Conf. Comput. Vis. Pattern Recognit. (CVPR)}, pp. 16739--16752, 2025.

\bibitem{c1} A. Sridhar, D. Shah, C. Glossop, and S. Levine, ``NoMaD: Goal Masked Diffusion Policies for Navigation and Exploration,'' in \emph{Proc. IEEE Int. Conf. Robot. Autom. (ICRA)}, 2024.

\bibitem{c2} Y. Deng, S. Yuan, and Y. Fang, ``AnyImageNav: Any-View Geometry for Precise Last-Meter Image-Goal Navigation,'' arXiv:2604.05351, 2026.

\bibitem{meng2025aim} X. Meng, X. Yang, S. Jung, F. Ramos, S. S. Jujjavarapu, S. Paul, and D. Fox, ``Aim My Robot: Precision Local Navigation to Any Object,'' in \emph{Proc. IEEE Int. Conf. Robot. Autom. (ICRA)}, 2025.

\bibitem{c3} C. Chi, Z. Xu, S. Feng, E. Cousineau, Y. Du, B. Burchfiel, R. Tedrake, and S. Song, ``Diffusion Policy: Visuomotor Policy Learning via Action Diffusion,'' in \emph{Proc. Robot Sci. Syst. (RSS)}, 2023.

\bibitem{jaegle2021perceiver} A. Jaegle, F. Gimeno, A. Brock, A. Zisserman, O. Vinyals, and J. Carreira, ``Perceiver: General Perception with Iterative Attention,'' in \emph{Proc. Int. Conf. Mach. Learn. (ICML)}, pp. 4651--4664, 2021.

\bibitem{peebles2023dit} W. Peebles and S. Xie, ``Scalable Diffusion Models with Transformers,'' in \emph{Proc. IEEE/CVF Int. Conf. Comput. Vis. (ICCV)}, pp. 4195--4205, 2023.

\bibitem{lu2022dpmsolverpp} C. Lu, Y. Zhou, F. Bao, J. Chen, C. Li, and J. Zhu, ``DPM-Solver++: Fast Solver for Guided Sampling of Diffusion Probabilistic Models,'' arXiv:2211.01095, 2022.

\bibitem{dino2025} O. Sim\'eoni \emph{et al.}, ``DINOv3,'' arXiv:2508.10104, 2025.

\end{thebibliography}

\end{document}
